%%%%%%%%%%%%%%%%%%%%%%%%%%%%%%%%%%%%%%%%%%%%%%%%%%%%%%%%%%%%%%%%%%%%%%%%%%%%%%%
% Globale Definitionen
%%%%%%%%%%%%%%%%%%%%%%%%%%%%%%%%%%%%%%%%%%%%%%%%%%%%%%%%%%%%%%%%%%%%%%%%%%%%%%%

%=== pdf Metadaten ============================================================
\hypersetup{
	pdfauthor={Kishanlal Rajesh Kanojia},
	pdftitle={Analysis of Service Tickets Using LLM},
	pdfsubject={Master Thesis},
	pdfkeywords={%
		Large Language Models,
		Service Tickets,
		Information Extraction,
		NLP,
		Ollama,
		Tag Classification,
	},
	bookmarksnumbered=true,
	pdfstartview=FitH,
	hidelinks,
}

%=== Vorwort vor Literatur ====================================================
\defbibnote{mynote}{%
	As explained in  \cref{sec:bib-content}, only the sources that are actually 
	referenced in a paper are listed in the bibliography. Therefore, please 
	always check the bibliography of your paper to ensure that all sources
	cited~--~and only these~--~have been included. This does not apply to the 
	following list; most of the sources listed here are not cited in this
	template. This is merely an example of a bibliography with recommended 
	reading for academic writing.
	
}

%=== Name für Inhaltsverzeichis ===============================================
\renewcaptionname{USenglish}{\contentsname}{Content}

%=== Kopf-/Fusszeile definieren ===============================================
\clearpairofpagestyles
\ohead[]{\headmark}
\ofoot[\pagemark]{\pagemark}

%=== Farben definieren ========================================================
\definecolor{THRed}{RGB}{207,24,32}
\definecolor{THOrange}{RGB}{236,101,37}
\definecolor{THPurple}{RGB}{175,54,140}

%=== Einstellungen für cref ===================================================
%\newcommand{\crefpairconjunction}{ und~}
%\newcommand{\crefrangeconjunction}{ bis~}


%=== Einstellungen für plots ==================================================
\pgfplotsset{
	compat=newest,
	/pgf/number format/.cd,
	dec sep={\text{,}},
	1000 sep={\,},
}

%=== Einstellungen für listings ===============================================
\lstdefinestyle{myLaTeX}{
	basicstyle=\footnotesize\ttfamily,
	language=TeX,
	keywordstyle=\color{blue},
	frame=single,
	backgroundcolor=\color{gray!10},
	tabsize=2,
	morekeywords={
		lstdefinestyle,
		footnotesize,
		ttfamily,
		color,
	},
}

\lstdefinestyle{myBasic}{
	basicstyle=\footnotesize\ttfamily,
	frame=single,
	escapechar={|_},
	backgroundcolor=\color{white},
	keywordstyle=\color{black},
}

\lstset{style=myLaTeX}

%\lstset{language=Python, basicstyle=\small, frame=single, numbers=left, xleftmargin=2em, framexleftmargin=1.5em}
%\renewcommand{\lstlistingname}{Algorithmus}


%=== paar convenience Sachen definieren =======================================
\DeclareMathOperator{\sgn}{sgn}
\newcommand{\vecW}{\ensuremath{\mathbf{w}}}%
\newcommand{\ci}{\ensuremath{\mathrm{i}}}%
%
\newcommand{\tb}{\textbackslash}
%
\newcommand{\comm}[1]{\enquote{\texttt{\tb #1}}}
%
\newcommand{\param}[1]{%
	$\langle$\textrm{\textit{#1}}$\rangle$%
}

%=== Arial als Hauptschriftart ================================================
%\setsansfont{Arial}
%\renewcommand{\familydefault}{\sfdefault}

%=== keine größeren Leerzeichen nach Punkt (bei englischer Spracheinstellung) ==
\frenchspacing
